\section{Charts and Comparisons}
\label{sec:charts-comparisons}

Figures~\ref{fig:runtime-outcomes} and~\ref{fig:metric-outcomes} summarize testable runtime and metric outcomes. The tables that follow report the environment ANOVA for several metric comparisons.
\subsection{Runtime by Data Structure}
\begin{figure}[H]
\centering
\begin{tikzpicture}
\begin{axis}[
  ieee stacked bars,
  bar width=26pt,
  xticklabels={Matrix, Uniform Nodes, Memory Stress, Linked List},
]
\addplot coordinates {(0,33.33333333) (1,8.00000000) (2,27.27272727) (3,0.00000000)};
\addlegendentry{Significant improvement}
\addplot coordinates {(0,23.33333333) (1,8.00000000) (2,0.00000000) (3,0.00000000)};
\addlegendentry{Non-significant improvement}
\addplot coordinates {(0,36.66666667) (1,84.00000000) (2,72.72727273) (3,100.00000000)};
\addlegendentry{Significant worsening}
\addplot coordinates {(0,6.66666667) (1,0.00000000) (2,0.00000000) (3,0.00000000)};
\addlegendentry{Non-significant worsening}
\end{axis}
\end{tikzpicture}
\caption{Runtime outcomes by data structure.}
\label{fig:runtime-outcomes}
\end{figure}

Figure~\ref{fig:runtime-outcomes} puts into perspective how effective the custom memory allocators are in improving the runtime of each of the workloads. 
Starting with the worst performance, the Linked Lists showed only significant worsening in all testing environments, run lengths, and allocator versions.
The Uniform Nodes workloads were second worst, with 84\% of the testable comparisons showing significant worsening, and 8\% showing non-significant improvement. Although it did improve in short and medium runs for version 3.0.0.
The Memory Stress workloads had about 27\% with significant improvement and the rest being significant worsening. Version 6.0.0 performed the best here as well, showing a consistent best showing from version 6.0.0.
Since the custom memory allocator is focused mainly on improving the performance of parallelized matrices, it had over 50\% improvement although about 20\% of that is non-significant.

\subsection{Metric Comparison}

\begin{figure}[H]
\centering
\begin{tikzpicture}
\begin{axis}[
  ieee stacked bars,
  bar width=18pt,
  enlarge x limits=0.10,
  x tick label style={font=\scriptsize, align=center, text width=1.7cm},
  xticklabels={Runtime, Peak Memory (KB), Major Faults, Minor Faults, Voluntary CS, Involuntary CS},
]
\addplot coordinates {(0,22.80701754) (1,0.00000000) (2,23.07692308) (3,70.83333333) (4,41.66666667) (5,54.16666667)};
\addlegendentry{Significant improvement}
\addplot coordinates {(0,14.03508772) (1,0.00000000) (2,0.00000000) (3,0.00000000) (4,0.00000000) (5,0.00000000)};
\addlegendentry{Non-significant improvement}
\addplot coordinates {(0,59.64912281) (1,100.00000000) (2,76.92307692) (3,29.16666667) (4,33.33333333) (5,45.83333333)};
\addlegendentry{Significant worsening}
\addplot coordinates {(0,3.50877193) (1,0.00000000) (2,0.00000000) (3,0.00000000) (4,25.00000000) (5,0.00000000)};
\addlegendentry{Other tested outcomes}
\end{axis}
\end{tikzpicture}
\caption{}
\label{fig:metric-outcomes}
\end{figure}

Figure~\ref{fig:metric-outcomes} focuses on additional measured metrics for the testing runs which focus on the weaknesses of certain workloads over each allocator. For example, the custom memory allocator always had a higher peak memory usage which only gets worse as the versions increased. This data however, is limited to team 2 tested versions 2, 7, and 8 although not enough data for a significance test was available with version 8.
One interesting metric from this is the difference between major and minor faults. While major faults increased significantly for the custom memory allocators, minor faults dropped up to 99+\% in version 7. However, for context switching, the custom memory allocator was better about preventing involuntary context switches while keeping it  about even in voluntary context switches.
In terms of runtime though, only about 20\% of runs had significant improvement.
For more data on the metrics see Table~\ref{PeakMemoryMetrics}, Table~\ref{MajorFaultsMetrics}, Table~\ref{MinorFaultsMetrics}, Table~\ref{VoluntaryContextSwitchMetrics}, and Table~\ref{InvoluntaryContextSwitchMetrics}.

\subsection{Environment Comparison Across Same Version and Workload}
The following table compares the same versions and workloads across multiple environments. However, this comes with the limitations that in order to reach the same run duration, the parameters may be different for some of these results.

\begin{table}[H]\caption{Environment ANOVA}
\centering\footnotesize\setlength{\tabcolsep}{5pt}\renewcommand{\arraystretch}{1.12}
\begin{tabular}{@{}llllcrrr@{}}
\toprule
Version & Workload & Comparison & Env.\ IDs & $n$ & $F$ & $p$ & Sig. \\ \midrule
v7.0.0 & Matrix & All Sizes & 1, 2 & 20 & 60.3979 & $3.6929\times10^{-7}$ & Yes \\
v7.0.0 & Matrix & Short & 1, 2 & 10 & 47.5496 & $1.2508\times10^{-4}$ & Yes \\
v6.0.0 & Matrix & All Sizes & 1, 2 & 20 & 374.4904 & $1.6999\times10^{-13}$ & Yes \\
v6.0.0 & Matrix & Short & 1, 2 & 10 & 111.9952 & $5.5525\times10^{-6}$ & Yes \\
v5.0.0 & Matrix & All Sizes & 1, 2 & 25 & 0.6923 & 0.4139 & No \\
v5.0.0 & Matrix & Short & 1, 2 & 10 & 0.5489 & 0.4799 & No \\
v5.0.0 & Matrix & Medium & 1, 2 & 10 & 0.4851 & 0.5058 & No \\
v5.0.0 & Uniform Nodes & All Sizes & 1, 2, 3 & 35 & 2,866.2357 & $8.1329\times10^{-37}$ & Yes \\
v5.0.0 & Uniform Nodes & Short & 1, 2, 3 & 15 & 874.0824 & $1.0041\times10^{-13}$ & Yes \\
v5.0.0 & Uniform Nodes & Medium & 1, 2 & 10 & 7,014.0642 & $4.6085\times10^{-13}$ & Yes \\
v5.0.0 & Uniform Nodes & Long & 1, 2 & 10 & 1,389.704 & $2.9414\times10^{-10}$ & Yes \\
v4.0.0 & Matrix & All Sizes & 1, 2 & 25 & 0.0879 & 0.7696 & No \\
v4.0.0 & Matrix & Short & 1, 2 & 10 & 35.7068 & $3.3235\times10^{-4}$ & Yes \\
v4.0.0 & Matrix & Medium & 1, 2 & 10 & 1.996 & 0.1954 & No \\
v4.0.0 & Uniform Nodes & All Sizes & 1, 2, 3 & 35 & 3,193.9375 & $1.4520\times10^{-37}$ & Yes \\
v4.0.0 & Uniform Nodes & Short & 1, 2, 3 & 15 & 1,376.1068 & $6.6936\times10^{-15}$ & Yes \\
v4.0.0 & Uniform Nodes & Medium & 1, 2 & 10 & 7,087.5386 & $4.4205\times10^{-13}$ & Yes \\
v4.0.0 & Uniform Nodes & Long & 1, 2 & 10 & 818.6869 & $2.4074\times10^{-9}$ & Yes \\
v3.0.0 & Matrix & All Sizes & 1, 2 & 15 & 1.452 & 0.2497 & No \\
v3.0.0 & Matrix & Short & 1, 2 & 10 & 4.1199 & 0.0769 & No \\
v2.0.0 & Matrix & All Sizes & 1, 2 & 20 & 4.6506 & 0.0448 & Yes \\
v2.0.0 & Matrix & Short & 1, 2 & 10 & 20.9244 & 0.0018 & Yes \\
\bottomrule\end{tabular}
\label{EnvAnova}
\end{table}

The key takeaways from Table~\ref{EnvAnova} are that all uniform nodes tests across environments had high levels of significance with p values much less than 10 to the 10th power. This means that there is a high likelyhood that the environment is playing a factor in the results.
However, the Matrix were split across the significance. This implies that across runs (mostly short and medium) there is little variation between systems which means that matrices are the better comparison metric when trying to determine across systems whether the custom memory allocator is effective or not.